\begin{theorem}[E.31 --- Proved: noncuspidal sections of the full profinite metabelian cubic quotient]
\label{E-full-metabelian-family}
Let $U=E\setminus E[3]$ be the plus Fermat cubic with origin
$O$, over $\mathbb Q$, and choose the rational tangential section
$s_O$ used above. In this statement $N_{\mathrm{all}}$ denotes
the full geometric profinite group, and
\[
 M=N_{\mathrm{all}}/N_{\mathrm{all}}'',\qquad
 A=N_{\mathrm{all}}^{\mathrm{ab}},\qquad
 D=N_{\mathrm{all}}'/N_{\mathrm{all}}''.
\]
There is an injective continuous family of arithmetic section classes
\[
 \mathbb Z_3\longrightarrow
 \operatorname{Sec}(M\rtimes_{s_O}G_{\mathbb Q})/M,
       \qquad a\longmapsto[s_a],
 \tag{E.31.1}
\]
whose abelianized section is $[s_O]$ for every $a$.
The class $[s_0]$ is $[s_O]$. Every $a\ne0$ gives a
noncuspidal section of this full profinite metabelian quotient.
Its images in every finite quotient, including quotients with
several coefficient primes, are the images of this one actual
continuous section.

The first nonzero family member is detected in the intrinsic
class-two three-primary quotient with coefficients modulo $27$.
The construction does not give a lift to the full geometric
fundamental group. In particular it supplies actual full
metabelian classes for task LM7-E.23(iii), while the complete
section image and its other fibres remain as stated in E.21.
\end{theorem}
\begin{proof}
We first construct the comparison that makes the full profinite
claim possible. Let $P=N_{\mathrm{all}}^{(3)}$ be the maximal
pro-three quotient and put
\[
 D_3=P'/P'',\qquad A_3=P^{\mathrm{ab}},\qquad
 R=\widehat{\mathbb Z}[[A]],\qquad R_3=\mathbb Z_3[[A_3]].
\]
The group $N_{\mathrm{all}}$ is free profinite of rank ten,
and $P$ is free pro-three of that rank. A chosen free basis
gives the continuous Alexander boundary sequences
\[
 \begin{aligned}
 0\longrightarrow D&\longrightarrow R^{10}
       \xrightarrow{\partial} I_A\longrightarrow0,\\
 0\longrightarrow D_3&\longrightarrow R_3^{10}
       \xrightarrow{\partial_3} I_{A_3}\longrightarrow0,
 \end{aligned}
 \qquad
 \partial(r_i)=\sum_i r_i(X_i-1).
 \tag{E.31.2}
\]
Here the ideals are the kernels of augmentation. One direct
verification uses the Cayley graphs of the finite abelian
covering quotients: their edge and vertex boundaries are the
displayed rows, and their first homology is the abelianization
of the covering group. Passing to all finite coefficients and
all such covers gives (E.31.2). To justify the limit on first
homology, a continuous finite character of $N_{\mathrm{all}}'$
has kernel containing $N_{\mathrm{all}}'\cap V$ for some open
normal subgroup $V$ of $N_{\mathrm{all}}$. It extends to
$N_{\mathrm{all}}'V$ by being trivial on $V$. This is
well-defined, and it is a homomorphism because
$[V,N_{\mathrm{all}}']$ is in that intersection. Thus these
abelian covers detect every finite character of the closed
commutator group. The same argument over pro-three quotients
gives the second sequence and their natural comparison.

In the coefficient-three factor $\mathbb Z_3[[A]]$, average
over the characteristic subgroup $A^{(3')}$, the product of
the other pro-prime factors. At each finite quotient its
order is prime to three, so its normalized Haar average
is a compatible idempotent. Extend that idempotent by zero
in the other coefficient factors of $R$ and call it $e$.
Then
\[
 eR\xrightarrow{\sim}R_3,
 \qquad eI_A\xrightarrow{\sim}I_{A_3},
 \qquad e(X_i-1)\longmapsto X_{i,3}-1.
 \tag{E.31.3}
\]
Multiplication by a central idempotent is exact.
Comparing the kernels of the two rows in (E.31.2)
therefore shows that the natural map $q:D\to D_3$
restricts to a topological isomorphism $eD\to D_3$.
Its inverse followed by inclusion is a specified continuous
section
\[
 j:D_3\lhook\joinrel\longrightarrow D,
                  \qquad qj=1.
 \tag{E.31.4}
\]
It is independent of the free basis: the map $q$ is natural,
and $e$ is preserved by every automorphism of $A$.
For the actual arithmetic action these same facts give
$g(e d)=e\,g(d)$ and $qg=gq$, hence $jg=gj$.
This is an equivariant embedding of the derived modules;
no section $P\to N_{\mathrm{all}}$ is asserted.

We now specify the cocycle. Use the actual arithmetic
quotient from E.16 and E.18,
\[
 \Gamma=P_{\mathrm{ar}}\rtimes\langle c\rangle,
 \qquad
 P_{\mathrm{ar}}=G_{\mathbb Q(\zeta_3),\{3\}}^{(3)}
                  =\langle\tau,\tau'\rangle_{\mathrm{pro}-3},
 \qquad c\tau c=\tau'.
 \tag{E.31.5}
\]
The action on the entire geometric pro-three group, hence
on $D_3$, factors through this quotient. The proof in E.16
checks the three-group residual action and the free arithmetic
source; it is not a replacement of an arbitrary arithmetic
action by its pro-three quotient.

Choose positive compatible roots for $\kappa=\kappa_3(3)$.
The Frobenius basis at seven can be chosen with
\[
                    \kappa(c)=0,\qquad
                    \kappa(\tau)=1,\qquad
                    \kappa(\tau')=-1.
 \tag{E.31.6}
\]
Here is the exact normalization. Start with the conjugate
Frobenii $\tau_0,\tau'_0$ at the primes above seven and put
$b=\kappa(\tau_0)$. The cubic residue computation
$3^2=2$ modulo seven, with $\zeta_3$ reducing to $2$,
gives $b\equiv1\pmod3$. Set
$r=\tau'_0\tau_0^{-1}$ and
$h=r^{(b^{-1}-1)/12}$. The exponent is in $\mathbb Z_3$.
One has $\chi(r)=1$, $\kappa(r)=-2b$, and
\[
 \kappa(h\tau_0h^{-1})
    =b+12b\,(b^{-1}-1)/12=1.
\]
Take $\tau=h\tau_0h^{-1}$ and $\tau'=c\tau c$.
Their Frattini classes are unchanged, so they remain a
free basis. This also checks the normalization used in
the supplied LM7 tower return, E.49, by its actual map.
The image on $V_3$ over $\mathbb Q(\zeta_3)$ is abelian:
after the explicit rational splitting of E.18 it is the
sum of the compact complex-multiplication character and
the boundary cyclotomic character. Equality on the rational
representation implies equality on its stable integral
lattice. Thus conjugating $\tau_0$ by $h$ leaves its
displayed matrix on $V_3$ unchanged; the finite certificate
below uses this same normalized generator.

Let $u,v\in V_3=P^{\mathrm{ab}}$ be the two compact
lift vectors of E.18, so $w=u\wedge v$ is primitive,
arithmetic $g$ acts on $w$ by $\chi_3(g)$, and
$c(w)=-w$. Choose free geometric generators lifting
the stated basis of $V_3$. Ordered products with the
integer exponent vectors $u,v$ give elements
$\widetilde u,\widetilde v\in P$. Put
\[
                  d=[\widetilde u,\widetilde v]\in D_3.
\]
Its image under $D_3\twoheadrightarrow
\gamma_2P/\gamma_3P=\bigwedge^2V_3$ is exactly $w$.
The free pro-three property in (E.31.5) supplies a unique
continuous crossed homomorphism on $P_{\mathrm{ar}}$
with values
\[
                       b(\tau)=d,\qquad b(\tau')=c(d).
\]
Indeed these values give a homomorphism into the pro-three
semidirect product $D_3\rtimes P_{\mathrm{ar}}$ on the two
free generators. Taking $b(c)=0$ extends it to $\Gamma$,
because the two assignments respect $c\tau c=\tau'$ and
$c^2=1$. Its degree-two projection is precisely
$\kappa w$, by (E.31.6) and equality on these generators.
Inflate it to $G_{\mathbb Q}$.

For $a\in\mathbb Z_3$ define the actual full metabelian
section by
\[
                     s_a(g)=j(a\,b(g))s_O(g).
 \tag{E.31.7}
\]
The image of $j$ is in the abelian normal subgroup $D$.
Its equivariance and the cocycle identity prove directly
that (E.31.7) is a continuous homomorphism and a splitting.
It has zero abelian cocycle. Its image in the three-primary
class-two fibre is $a[\kappa w]$.

For completeness this latter class is nonzero with an exact
finite certificate. In the ordered wedge basis of E.18,
write $T=\bigwedge^2 A_\tau$ and $J=\bigwedge^2 C_V$.
After the normalization $b(c)=0$, coboundaries are exactly
the image of $(T-I)(I+J)$. The program
\texttt{full\_metabelian\_family.py} produces an integral
row $\lambda$ with
\[
 \lambda(T-I)(I+J)=0,
                   \qquad \lambda w=1872\equiv9\pmod{27}.
 \tag{E.31.8}
\]
The row, both $45$-dimensional matrices and the wedge
vector are included in its output. Thus the first member
is already nonzero with these finite coefficients.
For any $a\ne0$, the same row modulo
$3^{v_3(a)+3}$ detects $a[\kappa w]$.
Equivalently, the rational argument uses
$(T-I)^{-1}w=w/6$ and $Jw=-w$; this vector cannot be
a $J$-fixed coboundary parameter.

There is no further identification by conjugation in $M$.
A conjugator taking two zero-abelian sections to one another
has abelian coordinate in $H^0(G_{\mathbb Q},A)$, which is zero.
For each coefficient prime this follows from the compact
weight-one and boundary weight-two eigenvalues at a good
prime distinct from it; no eigenvalue is one. Thus the
conjugator belongs to $D$, and it changes $jb$ only by a
coboundary. Projecting with $q$ and then to degree two
proves injectivity in (E.31.1). Continuity follows from
continuous scalar multiplication and evaluation of (E.31.7);
the compact parameter space and the closed profinite
conjugacy relation also give the asserted section-class topology.

Finally consider a cusp packet in this quotient. A rational
cusp $x\ne O$ has nonzero compact Kummer class of $x-O$:
it is a nonzero three-torsion point of $E(\mathbb Q)$,
and the Mordell--Weil Kummer map into $H^1(\mathbb Q,TE)$
is injective. Its packet therefore contains no zero-abelian
class. A nonrational cusp has no packet over $\mathbb Q$.
At $O$, the cusp line $I_O^{\mathrm{ab}}\simeq
\widehat{\mathbb Z}(1)$ is primitive in $A$, and
$H^0(G_{\mathbb Q},A/I_O^{\mathrm{ab}})=0$ by the same
weight calculation. The exact cohomology sequence makes
the packet map injective on abelianization. Its only
zero-abelian class is therefore the parameter-zero
section $s_O$. Equation (E.31.8) excludes that class
for every $a\ne0$. This proves noncuspidality in the
stated quotient and all assertions of the theorem.
The primitive-loop centralizer calculation in E.33
below verifies directly that the cusp normalizers
of this metabelian quotient have precisely these
packets; in particular no extra normalizer section
is omitted by the comparison with the full group.
The anchors are \S51.16, folios 733v--736v, scans172--178,
and task LM7-E.23(iii).
\end{proof}
